% 2017_national_paper_1_liberal_q10_fig1（与理科 q8 同构）。
% 节点清单：开始 / 输入 n=0 / A=3^n-2^n / 空白判断菱形(原图空白) /
%   空白处理框(原图空白) / 输出 n / 结束。
% 有向边：开始→输入；输入→J；J→A；A→菱形；菱形否(右)→输出；输出→结束；
%   菱形是(左)→空白框；空白框→J 回环。是/否位置与原图一致。
% 首轮后修复：空白菱形加 minimum height=1.5cm（空内容 diamond 收扁，
%   与原图完整菱形不符）；其余不动。
\begin{tikzpicture}[x=1cm,y=1cm]
\tikzset{
  flownode/.style={draw, line width=0.6pt, align=center, inner sep=3pt,
    execute at begin node={\setlength{\parindent}{0pt}}},
  st/.style={flownode, shape=rectangle, rounded corners=1.5mm, minimum width=26mm, minimum height=10mm, inner xsep=3mm, inner ysep=1.5mm},
  io/.style={flownode, trapezium, trapezium left angle=68, trapezium right angle=112,
    minimum width=3.0cm, minimum height=1.0cm},
  pr/.style={flownode, minimum width=3.0cm, minimum height=1.0cm},
  dc/.style={flownode, diamond, aspect=2, minimum width=3.4cm, inner sep=2pt},
  fe/.style={draw, line width=0.6pt, -{Stealth[length=2.8mm,width=2.2mm]}},
  fl/.style={draw, line width=0.6pt}
}
\node[st] (st0) at (0,0) {开始};
\node[io,minimum width=3.4cm] (io0) at (0,-1.75) {输入 $n=0$};
\coordinate (J) at (0,-2.95);
\node[pr,minimum width=3.8cm] (prA) at (0,-4.0) {$A=3^{n}-2^{n}$};
\node[dc,minimum width=5.0cm,minimum height=1.5cm,aspect=2.8] (dc0) at (0,-5.9) {};
\node[pr,minimum width=3.0cm] (prB) at (-4.3,-4.0) {};
\node[io,minimum width=2.8cm] (io1) at (0,-8.05) {输出 $n$};
\node[st] (st1) at (0,-9.55) {结束};
\draw[fe] (st0.south) -- (io0.north);
\draw[fe] (io0.south) -- (J) -- (prA.north);

\draw[fe] (prA.south) -- (dc0.north);
\draw[fe] (dc0.south) -- (io1.north);
\node[right=1.5mm,xshift=-0.4813mm,yshift=0.0631mm] at (0,-7.17) {否};
\draw[fe] (io1.south) -- (st1.north);
\draw[fe] (dc0.west) -- (-4.3,-5.9) -- (prB.south);
\node[above=2mm,xshift=-0.1055mm,yshift=-0.7318mm] at (-3.0,-5.9) {是};
\draw[fe] (prB.north) -- (-4.3,-2.95) -- (J);
\end{tikzpicture}
